\subsection{The unique residual pattern on three cross arrows}

\begin{lemma}[Commuting images on three points]
\label{C:three-point-commuting-images}
Let two finite transformation monoids act on a three-element set,
and suppose their actions commute.  At least one of the following
holds: they have a common fixed point; both act by permutations;
one acts trivially; there is a proper idempotent transformation
commuting with both; or, after relabeling the points and interchanging
the sides, their images are exactly
\[
 \{\id,t_0,t_1\},\qquad \{\id,c_0,c_1\},
\]
where $t_i$ fixes $0,1$ and sends $2$ to $i$, and $c_i$ is
constant with value $i$.
\end{lemma}
\begin{proof}
Here is a complete finite verification with two independent
enumerations.  Encode each of the $27$ transformations by its
three values.  The first enumeration starts with the identity
submonoid, adjoins each omitted transformation, takes product
closure, and repeats until no new closed subset occurs.  It gives
exactly $699$ submonoids.  Every submonoid occurs, since adjoining
its elements one at a time gives a path to it in this search.

The independent enumeration assigns each of the $27$ maps to
an included or excluded set.  It propagates the three implications
\[
 a,b\in P\Rightarrow ab\in P,\qquad
 a\in P,\ ab\notin P\Rightarrow b\notin P,\qquad
 b\in P,\ ab\notin P\Rightarrow a\notin P.
\]
It branches on every still undecided map, beginning with the
identity included.  Its $1{,}463$ branch nodes give exactly the
same $699$ literal subsets.  Their $1{,}854$ commuting ordered
pairs are tested directly.  In the displayed order of the
alternatives, their disjoint counts are
\[
\begin{array}{c|rrrrr}
\text{alternative}&\text{fixed point}&\text{permutations}&
 \text{trivial side}&\text{proper idempotent}&\text{residual}\\\hline
\text{ordered pairs}&595&5&1158&90&6.
\end{array}
\]
For the fourth test, all $27$ maps are checked for idempotence,
proper image, and literal commutation with both actions.  Its
ninety witnesses are retained.  The six remaining ordered pairs
are exactly the simultaneous state conjugates and side interchanges
of the pair in the statement.
The complete lists and verification are
\path{code/outputs/b_three_state_commuting_images.json} and
\path{code/outputs/b_three_state_image_classification_verified.json};
the two algorithms are
\path{code/b_three_state_commuting_images.py} and
\path{code/b_verify_three_state_image_classification.cpp}.
This classifies the action images only; it does not replace either
original corner monoid by its image.
\end{proof}

\begin{lemma}[The residual three-state collage has zero second homology]
\label{C:three-state-residual-H2zero}
Let $M,N$ be arbitrary finite monoids.  Suppose their commuting
actions on $S=\{0,1,2\}$ have right $M$-image
$\{\id,c_0,c_1\}$ and left $N$-image
$\{\id,t_0,t_1\}$, with notation as in
Lemma~\ref{C:three-point-commuting-images}.
Then the collage $\mathcal C_S$ has
$H_2(B\mathcal C_S;\mathbb Q)=0$.
\end{lemma}
\begin{proof}
Put $U=\{0,1\}$, $u=[1]-[0]$, and $v=[2]-[0]$.
The two nontrivial zero-one characters are
\[
 \chi(m)=\begin{cases}1& m\text{ acts identically},\\0&
 m\text{ acts constantly},\end{cases}
 \qquad
 \psi(n)=\begin{cases}1& n\text{ acts identically},\\0&
 n\text{ acts as }t_0\text{ or }t_1.\end{cases}
\]
As a right $M$-module, both $u$ and $v$ transform by $\chi$.
As a left $N$-module, $u$ is fixed and
\[
 n v=\psi(n)v+\gamma(n)u,
 \qquad \gamma(t_0)=0,\quad\gamma(t_1)=1.
\]
Here the last equalities mean any preimages of the indicated
transformations; $\gamma$ is zero on the identity image.
Let $F_N$ denote this two-dimensional left module.  Its extension
\begin{equation}\label{C:three-state-gamma-extension}
 0\longrightarrow\mathbb Q\longrightarrow F_N
 \longrightarrow\mathbb Q_\psi\longrightarrow0
\end{equation}
is not split.  Indeed a splitting would send the quotient generator
to $v+a u$.  A preimage of $t_0$ requires $a=0$, whereas a preimage
of $t_1$ requires $a=-1$.

Write $X_U\subseteq X_S$ for the classifying spaces of the
collages with cross-bisets $U\subseteq S$.
Their two vertex spaces are identical.  In the double action-bar
complex, the coefficient quotient is
\[
 \mathbb QS/\mathbb QU
   =\mathbb Q_\psi\otimes_{\mathbb Q}\mathbb Q_\chi.
\]
Consequently the reduced chains of the pair $(X_S,X_U)$ are
the suspension of the tensor product of the two bar complexes
representing
\[
 D_N=\mathbb Q\otimes^{\mathbf L}_{\mathbb QN}\mathbb Q_\psi,
 \qquad
 D_M=\mathbb Q_\chi\otimes^{\mathbf L}_{\mathbb QM}\mathbb Q.
\]
Both degree-zero homology groups vanish: each character takes
the value zero on some actor.  K\"unneth therefore gives
\[
 H_2(X_S,X_U;\mathbb Q)=0,
 \qquad
 H_3(X_S,X_U;\mathbb Q)
 =\operatorname{Tor}^{\mathbb QN}_1(\mathbb Q,\mathbb Q_\psi)
   \otimes
   \operatorname{Tor}^{\mathbb QM}_1(\mathbb Q_\chi,\mathbb Q).
\]

We identify the boundary of this pair explicitly.
The coefficient extension for $U\subseteq S$ is the pushout
of extensions of left-$N$, right-$M$ bimodules
\[
\begin{array}{ccccccccc}
0&\to&\mathbb Q\otimes\mathbb Q_\chi&\to&
 F_N\otimes\mathbb Q_\chi&\to&
 \mathbb Q_\psi\otimes\mathbb Q_\chi&\to&0\\
 &&\downarrow [1]-[0]&&\downarrow&&\Vert\\
0&\to&\mathbb QU&\to&\mathbb QS&\to&
 \mathbb Q_\psi\otimes\mathbb Q_\chi&\to&0.
\end{array}
\]
Thus the relative connecting map is the extension class
$\gamma$ of \eqref{C:three-state-gamma-extension}, tensored
with the identity on $\mathbb Q_\chi$, followed by the
inclusion $u:\mathbb Q_\chi\to\mathbb QU$.
This can also be read directly in normalized chains: in the
splitting $\mathbb QS=\mathbb QU\oplus\mathbb Qv$, the
right-$M$ faces have no off-diagonal term, and the only
off-diagonal left-$N$ face is $\gamma(n)u$.
Both vertex projections annihilate $v$.

On the core $U$, the left $N$-action is trivial.
The actual collapse $X_U\to B\mathcal C(M,U)$ is a rational
$H_2$ isomorphism by Lemma~\ref{C:trivial-biset-side-H2}.
The augmentation extension
\[
0\longrightarrow\mathbb Q_\chi
 \xrightarrow{u}\mathbb QU\longrightarrow\mathbb Q
 \longrightarrow0
\]
identifies the reduced act-cone complex with $D_M[1]$.
Its canonical quasi-isomorphism into the cone is induced by
$u$, not by a choice of homology splitting.  After the core
collapse, the boundary of the pair on degree three is therefore
\[
 \partial=\pm\gamma\otimes\id:
 \operatorname{Tor}^{\mathbb QN}_1(\mathbb Q,\mathbb Q_\psi)
 \otimes H_1(D_M)\longrightarrow H_1(D_M).
\]
The nonsplit extension gives a nonzero functional
$\gamma$ on the first Tor group, by the exact duality
\[
 \operatorname{Ext}^1_{\mathbb QN}(\mathbb Q_\psi,\mathbb Q)
 \cong
 \operatorname{Hom}_{\mathbb Q}\!\left(
   \operatorname{Tor}^{\mathbb QN}_1(\mathbb Q,\mathbb Q_\psi),
   \mathbb Q\right).
\]
That functional is onto $\mathbb Q$.  Thus $\partial$ is
onto $H_2(X_U;\mathbb Q)$, including when the latter is zero.
Together with $H_2(X_S,X_U;\mathbb Q)=0$, the homology
sequence of the pair proves $H_2(X_S;\mathbb Q)=0$.
\end{proof}

\begin{theorem}[Three cross arrows reduce to an ordinary act cone]
\label{C:three-cross-arrows-act-reduction}
If a two-object finite collage with three cross arrows has
type $\mathbb CP^2$, one side of the cross action is trivial,
and its corresponding ordinary three-state act cone also has
type $\mathbb CP^2$.
\end{theorem}
\begin{proof}
Apply Lemma~\ref{C:three-point-commuting-images} to the two
literal transformation images.  A common fixed state makes
the vertex degree-two groups inject into the edge; every
collage class is then relative and has zero square.
Permutation images are excluded by
Lemma~\ref{C:permutation-biset-cups}.
A proper equivariant idempotent retracts the collage to one
with at most two cross arrows; Lemma~\ref{C:proper-biset-retract-CP2}
would make that retract a projective plane, contrary to
Theorem~\ref{C:two-cross-arrows-universal} (the one-arrow
case is a join and has zero degree-two products).
The residual pattern is excluded by
Lemma~\ref{C:three-state-residual-H2zero}.
Only a trivial action side remains, and
Corollary~\ref{C:trivial-biset-side-CP2-retract} gives the
stated ordinary act cone.
\end{proof}

\begin{corollary}\label{C:three-cross-arrows-lower22}
Every finite projective-plane collage with three cross arrows
has at least twenty-two arrows.  If both corner monoids are
nontrivial, it has at least twenty-three arrows.
\end{corollary}
\begin{proof}
The ordinary three-state cone supplied by the theorem has at
least twenty-two arrows by Theorem~\ref{C:three-state-cone-minimum22}.
Removing a nontrivial trivially acting corner decreases the
number of arrows by at least one.
\end{proof}
